\section{Introduction and statement of the theorem}\label{sec:introduction}

The Boltzmann equation describes the leading empirical distribution of a
dilute gas. The next-order question asks whether the small correlations
created by deterministic collisions produce Gaussian fluctuations, and
which equation governs their covariance. We prove that the finite-dimensional
fluctuation law persists throughout an assumed regular kinetic interval,
without an equilibrium or small-data hypothesis.

Lanford's derivation of the Boltzmann equation established the kinetic
limit for short times~\cite{Lanford1975}. The next-order problem has a
separate history. Van Beijeren, Lanford, Lebowitz and Spohn derived
short-time equilibrium fluctuation covariances governed by the linearized
Boltzmann equation~\cite{vanBeijerenLanfordLebowitzSpohn1980}, and
Spohn established the nonequilibrium second-moment
limit~\cite{Spohn1981}. These covariance results identify the candidate
fluctuation dynamics; a Gaussian limit also requires control of
higher-order correlations.

Pulvirenti and Simonella developed short-time quantitative correlation-error
estimates and a cumulant-type reorganization of collision histories
\cite{PulvirentiSimonella2017}. Bodineau, Gallagher, Saint-Raymond and
Simonella developed the short-time nonequilibrium fluctuation theory
\cite{BGSS2020,BGSS2023}, including convergence of the fluctuation
process. Their connected-cumulant analysis identifies the Gaussian
collision noise and uses exact microscopic centering. Their direct
cluster expansion on physical trajectories~\cite{BGSSCluster2022}
provides a complementary representation of the short-time statistical
dynamics. At equilibrium, they obtained long-time covariance
convergence~\cite{BGSSCorrelations2023} and then the full Gaussian
fluctuation limit~\cite{BGSS2024}, using the invariant measure.
The BGSS fluctuation theorems concern a spatially periodic gas;
the present setting is the whole space, with integrable initial data.
Our conclusion extends the time interval away from equilibrium but
concerns finite-dimensional laws, rather than process convergence.

Deng, Hani and Ma extended the hard-sphere kinetic derivation through
each prescribed interval on which the regular Boltzmann solution
exists with the required bounds~\cite{DHM2025}. They also identify
long-time nonequilibrium fluctuations as a further application of
their approach~\cite[Section~1.4.3]{DHM2025}. The theorem below
establishes the finite-dimensional nonequilibrium fluctuation statement
on that regular kinetic lifespan. Its proof retains their layered
collision histories and supplies the sharper connected estimates and
centering transfer required at the fluctuation scale. It does not
assert global regularity for the Boltzmann equation or a path-space
tightness theorem.

\subsection{Kinetic equation and normalization}

Phase space is \(\R_x^3\times\R_v^3\), and \(z=(x,v)\).
For \(\beta>0\), define
\begin{equation}\label{eq:bol-norm}
 \|g\|_{\Bol,\beta}
 :=\sum_{a\in\Z^3}
   \sup_{\substack{|x-a|\le1\\v\in\R^3}}
    e^{\beta|v|^2}|g(x,v)|.
\end{equation}
The collision convention is
\begin{align}
 B(v-v_*,\omega)&=[(v-v_*)\cdot\omega]_+,
   &&\omega\in\mathbb S^2,\label{eq:collision-kernel}\\
 v'&=v-[(v-v_*)\cdot\omega]\omega,
 &v_*'&=v_*+[(v-v_*)\cdot\omega]\omega.\label{eq:scattering}
\end{align}
Here \(\dd\omega\) is the usual unnormalized surface measure, and
\([a]_+=\max\{a,0\}\). With all densities in the collision integral
evaluated at the same \(x\), set
\begin{equation}\label{eq:boltzmann}
 (\partial_t+v\cdot\nabla_x)f=Q(f,f),\qquad
 Q(f,f)=\int_{\R^3\times\mathbb S^2}
 B(v-v_*,\omega)(f'f_*'-ff_*)\,\dd v_*\,\dd\omega.
\end{equation}

\begin{assumption}\label{ass:kinetic}
The function \(f_0\ge0\) is a smooth probability density and, for some
\(\beta>0\),
\begin{equation}\label{eq:initial-bol-bounds}
 \|f_0\|_{\Bol,2\beta}
 +\|\nabla_xf_0\|_{\Bol,2\beta}<\infty.
\end{equation}
There is a nonnegative classical solution \(f\) of \eqref{eq:boltzmann} with initial
datum \(f_0\) on a prescribed finite interval \([0,T]\), \(T>0\), such that
\begin{equation}\label{eq:kinetic-gaussian}
 M:=\sup_{0\le t\le T}\sup_{x,v}
       e^{2\beta|v|^2}f(t,x,v)<\infty.
\end{equation}
\end{assumption}

\subsection{Microscopic law and fluctuation field}

For sphere diameter \(\eps>0\), let \(\mu=\eps^{-2}\), and write
\[
 \mathcal D_N^\eps
 =\{(z_1,\ldots,z_N):|x_i-x_j|>\eps\text{ for }i\ne j\}.
\]
The initial probability law, including \(N=0\), is
\begin{equation}\label{eq:grand-canonical}
 \P_\eps(N,\dd Z_N)
 =\mathcal Z_\eps^{-1}\frac{\mu^N}{N!}
   \ind_{\mathcal D_N^\eps}(Z_N)
   \prod_{i=1}^Nf_0(z_i)\,\dd z_i.
\end{equation}
The normalizing constant satisfies \(1\le\mathcal Z_\eps\le e^\mu\).
Conditional on the initial configuration, particles move by free flight
and deterministic elastic hard-sphere collisions. The almost-everywhere
definition of that flow is given by \cite[Proposition~1.2]{DHM2025};
its exceptional set has zero probability under the initial law.
There is no random collision rule
or resampling.

For real \(\phi\in C_c^\infty(\R^6)\), set
\begin{equation}\label{eq:empirical-fluctuation}
 \pi_t^\eps(\phi)=\mu^{-1}\sum_{i=1}^N\phi(z_i(t)),
 \qquad
 \zeta_t^\eps(\phi)
 =\sqrt\mu\,[\pi_t^\eps(\phi)-\E_\eps\pi_t^\eps(\phi)].
\end{equation}
The expectation in \eqref{eq:empirical-fluctuation} is exact. A
law-of-large-numbers approximation to it need not be accurate on this
scale; even the first excluded-volume correction at time zero can
survive after multiplication by \(\sqrt\mu\).

For an unmarked test \(\phi\), use the pair difference
\[
 \Delta\phi
 =\phi(x,v')+\phi(x,v_*')-\phi(x,v)-\phi(x,v_*).
\]
The limiting weak drift and collision covariance are
\begin{align}
 (A_t\phi)(x,v)
 &=v\cdot\nabla_x\phi(x,v)
   +\int_{\R^3\times\mathbb S^2}
      B(v-v_*,\omega)f_t(x,v_*)\Delta\phi
          \,\dd v_*\,\dd\omega,\label{eq:drift}\\
 C_t(\phi,\psi)
 &=\frac12\int_{\R^3\times\R^3\times\R^3\times\mathbb S^2}
       B(v-v_*,\omega)f_t(x,v)f_t(x,v_*)
       \Delta\phi\,\Delta\psi
       \,\dd x\,\dd v\,\dd v_*\,\dd\omega.\label{eq:noise-form}
\end{align}

\begin{theorem}[Finite-dimensional fluctuation limit]\label{thm:main}
Under \Cref{ass:kinetic}, the centered Gaussian generalized field
\(\zeta\) with weak equation
\begin{equation}\label{eq:fluctuating-equation}
 \dd\zeta_t(\phi)=\zeta_t(A_t\phi)\,\dd t+\dd W_t(\phi)
\end{equation}
is well defined by its backward mild form on a scale of tests with small
Gaussian growth. Its initial field is independent of \(W\), and
\begin{align}
 \E[\zeta_0(\phi)\zeta_0(\psi)]
   &=\int_{\R^6}f_0\phi\psi\,\dd z,\label{eq:initial-covariance}\\
 \E[W_t(\phi)W_s(\psi)]
   &=\int_0^{t\wedge s}C_r(\phi,\psi)\,\dd r.\label{eq:noise-covariance}
\end{align}
The Gaussian noise \(W\) has independent increments. For every integer
\(m\ge1\), fixed times \(t_1,\ldots,t_m\in[0,T]\), and real tests
\(\phi_1,\ldots,\phi_m\in C_c^\infty(\R^6)\),
\begin{equation}\label{eq:main-convergence}
 \bigl(\zeta_{t_1}^\eps(\phi_1),\ldots,
       \zeta_{t_m}^\eps(\phi_m)\bigr)
 \quad\Longrightarrow\quad
 \bigl(\zeta_{t_1}(\phi_1),\ldots,
       \zeta_{t_m}(\phi_m)\bigr)
 \qquad(\eps\downarrow0).
\end{equation}
\end{theorem}

The initial covariance in \eqref{eq:initial-covariance} has no subtraction
of \((\int f_0\phi)(\int f_0\psi)\): the particle count in
\eqref{eq:grand-canonical} is random. The equivalent backward covariance
formula and the precise construction of the Gaussian field are given in
\Cref{prop:gaussian-target}.

There is also a concrete distinction between microscopic and kinetic
centering. Put \(\varrho_0(x)=\int f_0(x,v)\,\dd v\).
The initial ensemble calculation in \Cref{prop:initial} gives
\[
 \lim_{\eps\to0}\sqrt\mu
 \left(\E\pi_0^\eps(\phi)-\int f_0\phi\right)
 =-\frac{4\pi}{3}\int f_0(x,v)\phi(x,v)\varrho_0(x)\,\dd x\,\dd v.
\]
Thus exclusion changes the mean at exactly the fluctuation scale,
even though its contribution to the limiting initial covariance
vanishes. Centering by the Boltzmann density alone would generally
introduce a nonzero deterministic shift already at time zero.

\subsection{Proof structure and new estimates}

The proof retains the layered collision-history framework of
\cite{DHM2025}; its operational inputs and their hypotheses are stated in
\Cref{thm:dhm-package}. The connected graph bounds also use the classical
tree-graph partition scheme of Penrose~\cite{Penrose1967}; we include its
proof for the initial ensemble. The fluctuation argument needs two
improvements beyond coarse decorrelation.

First, factorial cumulant measures on \(h\) observed particle histories
must be controlled at the sharp scale \(\mu^{1-h}\). We obtain an exact
component expansion that keeps finitely many observation records and
preserves factorial cancellation. After coarse bounds remove histories
of large complexity, chronological relative translations give one
\(\eps^2\) factor for every merger. Initial excluded-volume links gain
an additional \(\eps\), while a first redundant contact disappears on
each finite tree. A summable bound, uniform in the permitted history
sizes, justifies the passage from these finite histories to the full
expansion. This is \Cref{thm:sharp-cumulants}.

The surviving signed trees determine both one-particle means and the
connected two-particle contribution. Adding the diagonal term produces
\eqref{eq:noise-form}. Higher ordinary cumulants vanish after central-limit
scaling. A mark that accumulates observations along each particle reduces
every fixed multitime linear combination to this same calculation.
These steps first establish the Gaussian limit on the restarted,
truncated dynamics, without using a comparison with the true flow.

Second, the modified dynamics must approximate the true process strongly
enough to transfer its exact centering. We prove
\(\P_\eps(\text{any discrepancy before }T)=O(\eps^{1+\eta})\) for some
\(\eta>0\), in \Cref{prop:cutoff}. A primitive pair of close collisions
has a direct space--time--velocity gain. In the complementary case a
fixed number of improved local estimates dominates the free-component
costs. Joint energy and causal time-order bounds sum the ordinary
collision parameters without a logarithmic loss at every vertex.
These integration arguments apply to constrained histories without
assigning an independent Maxwellian to each cut variable.

Finally, this discrepancy estimate and fixed moments of \(N/\mu\)
transfer the limit and its exact centering to
\eqref{eq:empirical-fluctuation}. A probability of discrepancy tending
only to zero would suffice for uncentered weak comparison, but not for
this comparison of expectations on the fluctuation scale.

\Cref{sec:initial-target} treats the initial ensemble and constructs the
Gaussian target. \Cref{sec:dhm-inputs} states the layered-history
definitions and imported estimates, and \Cref{sec:cluster-expansion}
proves the marked component expansion. \Cref{sec:sharp-asymptotics}
derives the sharp connected bounds and tree limits, and
\Cref{sec:gaussian-limit} identifies the covariance and proves Gaussian
convergence for the truncated dynamics. \Cref{sec:cutoff} then proves
the discrepancy bound. \Cref{sec:transfer} transfers the exact centering
and completes the theorem for the true hard-sphere dynamics.
